%!TEX program = lualatex

\documentclass[11pt, a4paper]{article}
\usepackage[a4paper, margin=2.5cm]{geometry}
\usepackage{fontspec}
\usepackage[bulgarian, english, bidi=basic, provide=*]{babel}
\babelprovide[import, onchar=ids fonts]{bulgarian}
\babelfont{rm}{Noto Sans}
\babelfont{tt}{Noto Sans Mono}

\usepackage{xcolor}
\usepackage{hyperref}
\usepackage{booktabs}
\usepackage{enumitem}
\usepackage{listings}
\usepackage{titlesec}
\usepackage{xcolor}

\titleformat{\title}{\normalfont\Large\bfseries}{\thetitle}{1em}{}
\titleformat{\section}{\normalfont\large\bfseries}{\thesection}{1em}{}

\definecolor{primary}{RGB}{30, 64, 175}
\definecolor{secondary}{RGB}{71, 85, 105}
\definecolor{bgcode}{rgb}{0.93,0.93,0.93}
\definecolor{bordercode}{rgb}{0.82,0.82,0.82}

\hypersetup{
    colorlinks=true,
    linkcolor=primary,
    urlcolor=primary
}

\lstset{
    backgroundcolor=\color{bgcode},
    frame=single,
    rulecolor=\color{bordercode},
    basicstyle=\ttfamily\small,
    breaklines=true,
    showstringspaces=false
}

% Създаване на команда \code{...} за inline код със сив фон:
\newcommand{\code}[1]{\colorbox{bgcode}{\texttt{\detokenize{#1}}}}

\setlist[itemize]{label=-}

\title{\Large\bfseries Инструкции за инсталиране на Jupyter Notebook чрез Anaconda}
\date{}

\begin{document}

\maketitle

\section{Инсталация чрез Anaconda (препоръчително)}
\textbf{Anaconda} е цялостна платформа за Data Science, която включва Python, Jupyter Notebook и най-популярните библиотеки в един общ пакет.

\begin{enumerate}
    \item Отидете на официалния уебсайт: \url{https://www.anaconda.com/download}
    \item Регистрирайте се за да получите достъп.
    \item Изтеглете инсталационния файл \textbf{Anaconda distribution} за вашата операционна система (Windows, macOS или Linux).
    \item Стартирайте изтегления файл и следвайте инструкциите на екрана.
    \item След завършване на инсталацията вече разполагате с Python и Jupyter Notebook.
\end{enumerate}

\section{Стартиране на Jupyter Notebook}

Има няколко достъпни варианта за работа с Jupyter Notebook: локален \textbf{Jupyter сървър}, \textbf{VSCode}, \textbf{JupyterLab} или \textbf{Google Colab}.

\subsection{Локален Jupyter сървър}

През Anaconda Navigator:
\begin{enumerate}
    \item Отворете Anaconda Navigator.
    \item Кликнете върху иконата \textbf{Launch} под \textbf{Jupyter Notebook}.
    \item Това отваря Jupyter Notebook в папката на вашия потребителски профил.
    \item Създайте нова папка на удобно за вас място или използвайте вече съществуваща.
    \item Създайте нов ноутбук с \circ{\plus} бутона.
\end{enumerate}

\begin{enumerate}
    \item Отворете Терминал / Command Prompt.
    \item Отидете до папката, в която искате да работите. (напр. \code{cd ~/Documents})
    \item Изпълнете следната команда:
\end{enumerate}

\begin{lstlisting}
jupyter notebook
\end{lstlisting}

След изпълнение на командата, вашият подразбиращ се уеб браузър ще се отвори автоматично на адрес \url{http://localhost:8888}.

\subsubsection{Спиране на работата}
За да затворите Jupyter Notebook:
\begin{enumerate}
    \item Запазете файлове от менюто \textit{File} \rightarrow \ \ \textit{Save and Checkpoint}.
    \item В уеб браузъра натиснете бутона \textbf{Quit} в горния десен ъгъл.
    \item В терминала натиснете \texttt{Ctrl + C} и потвърдете с \texttt{y}, за да спрете сървъра.
\end{enumerate}

\subsection{За работа от VSCode}
\begin{enumerate}
    \item Отворете VSCode.
    \item Инсталирайте разширението \textbf{Python} от Microsoft, ако все още не е инсталирано.
    \item Отворете командния панел с \texttt{Ctrl + Shift + P} и изберете \textbf{Python: Select Interpreter}, след което изберете интерпретатора от Anaconda.
    \item Създайте нов файл с разширение \texttt{.ipynb} или отворете съществуващ Jupyter Notebook файл.
    \item Стартирайте клетките чрез бутоните за изпълнение или с \texttt{Shift + Enter}.
\end{enumerate}

\end{document}